% ---------------------------------------------------------------------------------------------------------------
% Discussion of the automatic-prompting results (Section 4.3). Requires: amsmath.
% ---------------------------------------------------------------------------------------------------------------

\section{Discussion of Automatic Prompting}
\label{sec:discussion_prompting}

\subsection{Why SAM Cannot Improve on Its Prompter}

The central finding of Section~\ref{sec:auto_prompt_results} is that SAM with automatic prompts is bounded by its
prompter. This follows from what the Phase-1 decoder was trained to do. It learned
$p(\text{mask}\mid\text{image},\text{prompt})$ from prompts that were always correct: every box contained exactly one
object, and the extent of the box was the extent of the object. Under this training distribution the optimal decoder
\emph{trusts} its prompt: it segments whatever lies inside a box and never has a reason to output nothing or to extend
a mask beyond its prompt. The prompt analysis confirms this behaviour directly (Figure~\ref{fig:oracle_prompt_errors}).
The errors of a prompter are of exactly these kinds, namely false objects, missed objects and wrong AR extents
(Figures~\ref{fig:object_level} and~\ref{fig:error_decomposition}), so they pass through SAM unchanged. What SAM can add
is boundary quality inside a correct prompt; but for ARs the dominant error is the extent of the whole river, not its
boundary pixels, and for TCs the objects are small enough that detection dominates.

The decoder-adaptation experiments show that this is not only a question of the training distribution. Even a decoder
trained on realistic, out-of-fold prompts learns to reject some false prompts but cannot add information that the
prompt does not contain. Prompter and decoder read the same frozen image embedding; whatever the decoder could infer
about the extent of an AR from that embedding, a prompter trained directly on the labels infers as well.

A further reason lies in the labels. ClimateNet's labels are the consensus of several experts, and the mean performance
of a single expert against this consensus is AR~0.341 / TC~0.257 IoU (Table~2.1). All good prompters of this study
(TC 0.33--0.35, AR 0.38--0.41) already exceed the average single expert. The remaining errors, such as the extent of a
river, where it starts and ends, or weak cyclones, are precisely where experts disagree, so part of the gap to the
oracle is irreducible label ambiguity rather than a deficiency of the model.

\subsection{Why a Small Generator Beats the Larger Ones}

The encoder describes the image on a $64\times64$ grid of $16\times16$-pixel patches. The ground truth has a resolution
of $768\times1152$, but its information content at object boundaries is limited by the ambiguity between experts. MSF
upsamples the features to $1024\times1024$ by nearest-neighbour interpolation and runs most of its 702~GFLOPs there,
which cannot create information and makes the model slow to train. MPG spends its capacity where the information is, at
$64\times64$, with a receptive field of about 15 patches from the dilated convolutions on top of the encoder's own global
attention, and uses two learned upsampling steps only to reach the $256\times256$ resolution of SAM's dense prompt. The
strong linear probe on the last ViT block supports this view: the adapted encoder has already done most of the work,
and the prompter mainly aggregates context and removes noise.

\subsection{Lessons for Prompt Design}

\begin{enumerate}
  \item If sparse prompts are used, they should be tight boxes derived from a segmentation. Points, loose boxes and
  boxes regressed by a detector are harmful.
  \item Dense prompts are more robust to prompt errors and carry the shape of an AR, but they must be passed as logits
  in the value range the decoder was trained with, and small TCs need an accompanying box on checkpoints whose decoder
  ignores weak dense evidence.
  \item One dense prompt per class is sufficient. This makes SAM decoding cheap: one decoder call per class and image
  instead of one per object.
\end{enumerate}

\subsection{Implications for a Prompt-free ClimateSAM}

For a prompt-free system, the prompter's own mask is the natural output (MPG: TC~0.344 / AR~0.413, mean IoU 0.567
including background). SAM with hybrid prompts gives masks of statistically similar quality ($-0.008$ mean FG IoU) with
smoother boundaries, at the cost of two decoder calls per image. The contribution of the foundation model is thus its
\emph{encoder}: the adapted features allow a 0.67\,M-parameter head to match a dedicated segmentation network trained
under the same protocol on the raw fields (mean FG IoU 0.379 versus 0.369--0.375, better for ARs by 0.013--0.018), and
a linear probe to come within about 0.03 of it. In terms of cost, however, CG-Net needs 23~GFLOPs per image, the frozen
encoder plus MPG about 980~GFLOPs. For a purely automatic system, a well-trained small network is therefore the cheaper
choice at nearly the same quality, and the two are complementary: their ensemble is the best automatic mask of this
study (0.389). The promptable decoder remains valuable for interactive use,
where a human annotator provides the boxes (Table~4.12: TC~0.72 / AR~0.63--0.64 with ground-truth boxes), but not as a
refinement stage after an automatic prompter.

\subsection{Limitations}

\begin{itemize}
  \item \emph{Small test set.} With 61 test images, differences below about 0.01 IoU lie within the bootstrap
  intervals and should not be interpreted. Checkpoints are selected on only 40 validation images drawn from the same
  years as the training images.
  \item \emph{Seeds.} Three seeds per configuration on the primary checkpoint; on the second checkpoint the MSF models
  were trained with one seed. The seed variance (standard deviation 0.001--0.01) is small compared with the variance
  from sampling test images.
  \item \emph{Two frozen checkpoints.} Both are Infused-Token ViT-B models; LoRA-adapted and ViT-L encoders were not
  tested in Phase~2.
  \item \emph{External prompter.} The fine-tuned CG-Net was trained on all 398 training images, including the
  validation images of this protocol; this affects only experiments that use its outputs on training images.
  \item \emph{Object metrics.} An object counts as found if one pixel overlaps; stricter matching (IoU $\ge 0.5$) would
  lower all recalls and precisions but not change the comparison.
  \item \emph{Earlier Table~4.13.} The re-run shows that its SAM rows do not reproduce with the current code and
  checkpoints; their origin could not be identified.
\end{itemize}

\subsection{Future Work}

\begin{enumerate}
  \item \emph{Train Phase~1 with realistic prompts.} Mixing ground-truth prompts with corrupted ones, including false
  boxes whose target is empty, and with prompts of an out-of-fold prompter already during Phase~1 attacks the root cause
  identified above. The prompt analysis of Section~\ref{sec:prompt_analysis} provides the test: a better decoder
  flattens the curves of Figure~\ref{fig:oracle_prompt_errors}.
  \item \emph{Query-based prompt-free decoding.} Replacing static prompts by image-dependent object queries
  (DETR/Mask2Former style) with a matching-based loss turns SAM into an instance segmenter and addresses the TC detection
  errors directly.
  \item \emph{Joint fine-tuning of encoder adapters and prompter.} Unfreezing the infused-token adapters together with
  MPG on the segmentation loss lets the encoder adapt to the prompt-free objective.
  \item \emph{Longitudinal periodicity.} Circular padding in the prompter and test-time averaging over longitudinally
  rolled inputs would fix objects cut at the date line.
  \item \emph{Stricter object metrics.} Object recall and precision at IoU $\ge 0.3$ / $0.5$ and event counts per image
  would make the detection analysis comparable with the ClimateNet literature.
\end{enumerate}
